\subsection[Radix Sort]{Radix Sort}
    \subsubsection{Overview}
        \begin{itemize}
            \item \textbf{History and Origin:} A non-comparative sorting algorithm originally used for sorting punch cards.
            \item \textbf{Efficiency:} Highly efficient for datasets with keys that are not too large (e.g., integers or strings). 
            \item \textbf{Application:} Widely used in environments where numerical data needs to be sorted in linear time relative to the number of digits.
        \end{itemize}
    \subsubsection{Idea}
        \begin{itemize}
            \item Radix Sort avoids direct comparison between elements. Instead, it groups elements by their individual digits.
            \item The implemented version is Least Significant Digit (LSD) Radix Sort, which processes digits from the least significant (rightmost) to the most significant (leftmost).
            \item It relies on a stable sorting subroutine (like Counting Sort) to sort the digits at each positional place value.
        \end{itemize}

    \subsubsection{Step-by-Step Algorithm}
        \begin{itemize}
            \item Consider an array of n elements.
            \item \textbf{Step 1:} Find the maximum value (\textit{maxv}) in the array to determine the number of digits.
            \item \textbf{Step 2:} Initialize a multiplier \textit{j} = 1.
            \item \textbf{Step 3:} While \textit{maxv} / \textit{j} > 0, perform Counting Sort for the current digit place:
            \begin{itemize}
                \item Create a frequency array \textit{Count} of size 10 (for base 10) initialized to 0.
                \item Calculate occurrences of each digit: \textit{digit} = (array[i] / j) \% 10.
                \item Accumulate \textit{Count} array to determine actual positions.
                \item Place elements into a temporary array \textit{B} traversing backwards to maintain stability.
                \item Copy array \textit{B} back to the original array.
            \end{itemize}
            \item \textbf{Step 4:} Multiply \textit{j} by 10 and return to \textbf{Step 3} until all digits are processed.
        \end{itemize}

    \subsubsection{Pseudo Code}
        \begin{itemize}
            \item \textbf{Algorithm:} LSD Radix Sort
            \item \textbf{Input:} An array A of n elements
        \end{itemize}
        \begin{verbatim}
        1. Find maxv in A
        2. j = 1
        3. While (maxv / j) > 0:
        4.    Count = array of 10 zeros
        5.    For i = 0 to n - 1:
        6.       digit = (A[i] / j) % 10
        7.       Count[digit]++
        8.    For k = 1 to 9:
        9.       Count[k] = Count[k] + Count[k - 1]
        10.   B = array of size n
        11.   For i = n - 1 down to 0:
        12.      digit = (A[i] / j) % 10
        13.      B[Count[digit] - 1] = A[i]
        14.      Count[digit]--
        15.   Copy B back to A
        16.   j = j * 10
        \end{verbatim}
    \subsubsection{Complexity}
        \begin{itemize}
            \item \textbf{Time Complexity:}
            \begin{itemize}
                \item Best, Average, and Worst case: $O(d(n + k))$
                    \subitem Where \textit{d} is the number of digits, \textit{n} is the number of elements, and \textit{k} is the base (10).
            \end{itemize}
            \item \textbf{Space Complexity:} $O(n + k)$ 
                \subitem Requires auxiliary arrays (\textit{Count} and \textit{B}) for the internal Counting Sort.
        \end{itemize}
    \subsubsection{Variants and Optimizations}
        \begin{itemize}
            \item \textbf{Current Implementation (Base-10 LSD):}
            \begin{itemize}
                \item The provided C++ code implements Least Significant Digit (LSD) Radix Sort using base 10 (decimal system). It extracts digits using modulo (\texttt{\% 10}) and allocates a counting array of size 10 for each pass.
            \end{itemize}
            \item \textbf{Theoretical Variant: Bitwise Radix Sort (Base-256):}
            \begin{itemize}
                \item \textbf{Problem:} Using base 10 requires division (\texttt{/}) and modulo (\texttt{\%}) operations, which are computationally expensive. Additionally, base 10 requires more passes for large integer limits.
                \item \textbf{Solution:} Change the radix base to a power of 2 (e.g., base 256 for byte-level extraction) using fast bitwise operations (right shift \texttt{>>} and bitwise AND \texttt{\&}).
                \item \textbf{Benefit:} Drastically reduces the number of sorting passes (e.g., only 4 passes for 32-bit integers) and replaces slow arithmetic operations with fast bitwise logic.
            \end{itemize}
        \end{itemize}